\begin{textsample}{NEG Dim 6 – human – Score -2.00 – mg/mg\_ef\_ciencias\_3\_p4.txt}  \label{ex:f6_neg_020}
O ensino de Ciências da Natureza no Ensino Fundamental, tem como intencionalidade cooperar para a transformação da sociedade, ao tratar dos saberes que lhes são inerentes, permitindo aos estudantes o desenvolvimento de habilidades para a construção, reconstrução ou desconstrução dos conhecimentos, premissa que requer a implementação de um conjunto de encaminhamentos que contribuam para a formação de estudantes questionadores e investigativos.
Essa capacidade investigativa dos estudantes, por sua vez, se dá por meio da observação e da pesquisa com o objetivo da integração do conhecimento científico resultante da investigação da natureza para interpretar racionalmente os fenômenos naturais \textbf{observados}, decorrentes de contextos históricos, sociais e econômicos considerados no tempo, espaço, matéria, movimento, força, campo, energia, vida e evolução.

% matched lemmas: observar
\end{textsample}
